% Standalone uncompiled research note. No external TeX inputs are required.
\documentclass[11pt]{article}
\usepackage[T1]{fontenc}
\usepackage[utf8]{inputenc}
\usepackage[margin=1in]{geometry}
\usepackage{amsmath,amssymb,amsthm,mathtools,enumitem}
\usepackage[hidelinks]{hyperref}
\newtheorem{theorem}{Theorem}[section]
\newtheorem{proposition}[theorem]{Proposition}
\newtheorem{corollary}[theorem]{Corollary}
\newtheorem{lemma}[theorem]{Lemma}
\theoremstyle{definition}\newtheorem{definition}[theorem]{Definition}
\newtheorem{example}[theorem]{Example}
\theoremstyle{remark}\newtheorem{remark}[theorem]{Remark}
\newcommand{\Pow}{\mathcal P}
\newcommand{\Gen}{\operatorname{Gen}}
\newcommand{\Meet}{\operatorname{Meet}}
\newcommand{\cl}{\operatorname{cl}}
\newcommand{\im}{\operatorname{im}}
\newcommand{\Set}{\mathbf{Set}}
\newcommand{\Cat}{\mathbf{Cat}}
\newcommand{\Bool}{\mathbf{2}}
\newcommand{\N}{\mathbb N}
\newcommand{\NN}{\mathbb N}
\newcommand{\Cost}{\mathbb N\cup\{\infty\}}
\newcommand{\Lan}{\operatorname{Lan}}
\newcommand{\Fact}{\operatorname{Fact}}
\newcommand{\id}{\mathrm{id}}
\title{Changing experiments and composing characterizations}
\author{A research note with exact criteria and native witnesses}
\date{5 November 2026}
\begin{document}
\maketitle
\begin{abstract}
Changing the experiments or operations available to a system can alter its
behavioral quotient, concept completion, and incidence topology in different
directions and at different times. We derive these changes from native
continuation tests, prove the least simultaneous operational repair, and
characterize when repair commutes with an actual wiring. A two-bit gate
example leaves the state quotient fixed while forcing a larger concept
lattice; a subsequent native reset changes only the witness topology.
For experiment formation by declared decompositions, the comparison between
changing interfaces and composing is a pointwise Kan-extension mate.
Its established exact-square criterion leads to a complete classification
for integer resource translations: Boolean feasibility, set-valued exactness,
and equivalence of retained allocation categories require successively
stronger matrix conditions. A response-preserving cone states the precise
bridge from these experiment comparisons to the heterogeneous invariants.
Finally, synchronized weighted native machines show that complete Boolean
agreement can hide an unbounded resource defect; a common endpoint-preserving
primitive normalization gives an effective witness compiler for arbitrary
finite runs. General categorical machinery and previous project results are
credited explicitly. The concrete classifications and their integration are
proved here without a claim of historical priority.
\end{abstract}

\section{The question and the level of the answer}\label{sec:question}
An experimental characterization is relative to what can be done and what
can be observed. Suppose an old repertoire is enlarged by a gate, an
instrument, a continuation, or a changed interface. Does the old
characterization still support the new operations? If components are then
wired, does changing their characterizations first give the same result as
changing the complete system? The answer depends on both the composition
and the information retained in the characterization.

The first part of this note uses actual deterministic programs and their
Boolean responses. It derives an enlarged test doctrine, rather than
postulating a refinement map, and compares predictive states, formal
concepts and supported topology. The second part studies experiment
formation by a supplied category of decompositions. Here the relevant
Beck--Chevalley comparison can fail through an absent lift, disconnected
lifts, or a noninvertible witness that an extensional colimit forgets.
The third part makes native implementation explicit in finite weighted
machines. It separates independently optimal local implementations from
one jointly executable implementation and supplies a finite primitive
certificate for their agreement.

These are different presentations of a common problem, not interchangeable
names for one construction. In particular, FCA completion is not asserted
to be Day convolution, and CNOT is not assumed to be a Kan-extension
composition. The bridge is an evaluation map from generated experiment
labels to their actual response columns; its hypotheses are proved or
stated where it is used.

\paragraph{Reading route.}
Sections~\ref{sec:source}--\ref{sec:cnot} derive the mixed changes and a
fully computed interaction. Sections~\ref{sec:decomposition}--\ref{sec:slice}
give the categorical comparison, resource classification and response
bridge. Section~\ref{sec:native} supplies native witness conditions.
Section~\ref{sec:contextads} gives the exact contextad connection, and
Sections~\ref{sec:prior}--\ref{sec:limits} delimit the increment and the
remaining question. The source is deliberately self-contained and
uncompiled. Finite checks and intermediate research files accompany it.


\section{Native programs generate the experimental doctrine}\label{sec:source}
The question is not merely how three invariants of one table differ.
An interaction can make a new experiment executable, and a new operation
can turn an old observation into a new experiment by precomposition.
We derive the resulting comparisons from these source operations. There
are two comparison directions and an essential distinction between
strict, lax, and partial transport.

Let $\mathcal C$ be a small category of admitted native programs and
$X:\mathcal C\to\mathsf{Set}$ their deterministic state semantics.
Objects are experimental interfaces. At interface $i$, supply a set
$O_i$ of labelled seed observations, with their actual Boolean response
columns $C_e\subseteq X_i$. A label includes a native observation
witness; a morphism of $\mathcal C$ includes its native program witness.
The functor $X$ need not be faithful: two different native programs may
have the same state transformation.

This is a relative source model, not a theorem that finitely many
observations determine all physical laws or all possible programs.
Actual residual stock, damaged specimens, and installed devices belong
in the source state when they affect subsequent execution. A partial
command can be represented on a \emph{certified domain interface} on
which it is total; this does not turn unsafe disabledness into an
observable outcome. A measurement may have a state-changing native
implementation even though its final Boolean response alone supplies
an incidence column. In particular, intersections of columns below are
semantic conjunctions, not automatically jointly executable tests on
one destructive specimen.

\begin{definition}[Generated experiments]
A generated experiment at $i$ is a labelled pair $(p,e)$ with
$p:i\to j$ and $e\in O_j$, evaluated by
\[
 C_{(p,e)}=X(p)^{-1}(C_e).
\]
Write $G_i$ for the family of its distinct response columns, while
retaining the pairs as native witness annotations. Define
\[
 M_i=\Meet(G_i)
 =\left\{\bigcap\mathcal U:\mathcal U\subseteq G_i\right\},
 \qquad
 \cl_i(A)=\bigcap\{C\in G_i:A\subseteq C\}.
\]
Empty intersections are $X_i$. Define $x\sim_i y$ when all columns of
$G_i$ have the same membership value at $x$ and $y$, and put
$Q_i=X_i/{\sim_i}$.
\end{definition}
The family $G$ is the least family containing the seeds and stable
under preimage along admitted programs. Precomposition acts on
\emph{labels} by $(q,e)\mapsto(qp,e)$, so the semantic construction has
actual native witnesses, not selected endpoint arrows.

\section{Generated state and concept actions}
\begin{theorem}[The source-derived operational doctrine]
For every $p:i\to j$:
\begin{enumerate}
\item $X(p)$ descends uniquely to $\bar p:Q_i\to Q_j$. These maps
make $Q$ a functor, with quotient maps $q_i:X_i\twoheadrightarrow Q_i$
natural in $i$.
\item Preimage restricts to a complete-meet map
$p^*:M_j\to M_i$. Its left adjoint is
\[
 p_!(A)=\cl_j(X(p)[A]).
\]
It preserves all joins and satisfies
$p_!\gamma_i(x)=\gamma_j(X(p)x)$, where
$\gamma_i(x)=\cl_i(\{x\})$.
\item The maps $p_!$ preserve identities and composition. They are the
unique complete-join actions agreeing with native state operations on
all marked object concepts $\gamma_i(x)$.
\end{enumerate}
\end{theorem}
\begin{proof}
A target generated column is $X(r)^{-1}(C_e)$, whose preimage under
$X(p)$ is $X(rp)^{-1}(C_e)\in G_i$. Consequently source-row equality
implies target-row equality. Surjectivity of $q_i$ gives uniqueness and
functoriality of the quotient operations.

Every extent is an intersection of generated columns; preimage preserves
arbitrary intersections, including the empty intersection. It therefore
restricts as claimed. For extents $A\in M_i,B\in M_j$,
\[
 \cl_j(X(p)[A])\subseteq B
 \quad\Longleftrightarrow\quad
 A\subseteq X(p)^{-1}(B),
\]
which proves the adjunction. Since $X(p)^{-1}(\cl_j(X(p)[U]))$ is a
closed set containing $U$, it contains $\cl_i(U)$. Thus
\begin{equation}
 \cl_j(X(p)[\cl_i(U)])=\cl_j(X(p)[U]).\tag{CN}
\end{equation}
This yields the marking equation and the composition law.
Finally $A=\bigvee_{x\in A}\gamma_i(x)$, with joins given by closed
unions. Every complete-join map agreeing on these generators is forced.
\end{proof}

For a one-object category this is a deterministic action of a monoid,
and generated tests are continuations. For a DFA with accepting set $F$,
$G=\{u_w^{-1}(F):w\in\Sigma^*\}$. The quotient is the predictive
row quotient; after retaining only reachable states it is the usual
minimal DFA. The lattice is its FCA/complete-meet completion, not the
state set itself. These fixed-interface facts are established automata,
FCA and Chu-space theory and were already proved in the preceding
\emph{common generation} manuscript.

\section{Growth forces a mixed variance comparison}\label{sec:growth}
Let $J:\mathcal C\hookrightarrow\mathcal D$ extend the admitted
program category, possibly by adding interface objects as well as
programs. Let $Y:\mathcal D\to\mathsf{Set}$ extend the source semantics,
with $YJ=X$. Old seed observations are retained at old interfaces;
new seeds may be supplied at any interface of $\mathcal D$. Decorate
new generated columns, lattices, closures and quotients by $+$.
At every old interface,
\[
 G_i\subseteq G_i^+,\qquad M_i\subseteq M_i^+.
\]
The new experiments can pass through newly introduced interfaces.
They need not be old experiments with renamed labels.

\begin{theorem}[Two directions, and strict backward coherence]
At each old interface the extension induces:
\begin{enumerate}
\item a unique surjection $\rho_i^Q:Q_i^+\twoheadrightarrow Q_i$
with $\rho_i^Qq_i^+=q_i$;
\item a complete-meet inclusion $j_i:M_i\hookrightarrow M_i^+$,
whose left adjoint is the surjective complete-join map
\[
 \rho_i^M:M_i^+\twoheadrightarrow M_i,
 \qquad \rho_i^M(A)=\cl_i(A);
\]
\item strict naturality for every \emph{shared} native program $p:i\to j$:
\begin{align*}
 \rho_j^Q\bar p^+&=\bar p\rho_i^Q,\\
 \rho_j^M p_!^+&=p_!\rho_i^M,\\
 \rho_i^M\gamma_i^+(x)&=\gamma_i(x);
\end{align*}
\item forward lax naturality, in the opposite direction,
\begin{align*}
 p_!^+j_i&\leq j_jp_!,\\
 \gamma_i^+(x)&\subseteq j_i\gamma_i(x).
\end{align*}
All backward comparisons compose strictly under successive extensions.
The forward inequalities are compatible with identities, composition
of shared programs, and successive extensions.
\end{enumerate}
\end{theorem}
\begin{proof}
Equality on more tests implies equality on fewer, proving (1).
For $A\in M_i^+$ and $B\in M_i$,
$\cl_i(A)\subseteq B$ iff $A\subseteq B$, so
$\rho_i^M\dashv j_i$. It is surjective because it fixes $M_i$.
The state square follows by evaluating it on all $x\in X_i$.
On right adjoints there is an equality
$j_i p^*=p^{*+}j_j$, since both sides are the same set-theoretic
inverse image. Taking left adjoints gives the second square. The
marking square follows from
$\cl_i\cl_i^+=\cl_i$.

New closure is smaller: $\cl_j^+(U)\subseteq\cl_j(U)$.
Apply this to $U=X(p)[A]$ to obtain the forward action inequality;
the marking inequality is the singleton case. For two shared programs,
monotonicity gives
\[
 q_!^+p_!^+j_i\leq q_!^+j_jp_!
                  \leq j_kq_!p_!.
\]
These are precisely the unique order comparisons between their
specified endpoints, so they respect the strict program composition
law. For successive growth stages $0\subseteq1\subseteq2$,
$\cl_0\cl_1=\cl_0$ and the inclusions compose literally;
the quotient comparison is unique on source-state representatives.
This proves all stated coherences.
\end{proof}
A newly added program need not act on the old quotient or lattice at all.
The theorem asserts naturality for shared operations, not an impossible
old action for every newly executable operation.

\begin{theorem}[Least simultaneous operational repair]
The generated extent family is the least Moore family containing its
seed columns for which every admitted source program has a
complete-join action agreeing with that program on marked object
concepts. Relative to an old interface, it is the least such extension
of the old extent family and all new seed columns. Independently, its
row equivalence is the largest program-stable equivalence refining the
seed distinctions, hence the coarsest adequate deterministic quotient.
\end{theorem}
\begin{proof}
For any Moore families $N_i$, a source map $f:X_i\to X_j$ has a
marking-compatible complete-join extension iff
\[
 f^{-1}(B)\in N_i\quad\text{for every }B\in N_j.\tag{E}
\]
For necessity, a join extension implies
$f[\cl_{N_i}(U)]\subseteq\cl_{N_j}(f[U])$, because
$\cl_{N_i}(U)$ is the join of the generators of $U$. Set
$U=f^{-1}(B)$: its closure still maps into $B$, so $U$ is closed.
For sufficiency, condition (E) proves (CN) and the adjunction exactly
as above. If every native primitive satisfies (E), induction on
programs forces every generated seed preimage into $N_i$; closure
under intersections forces $M_i\subseteq N_i$. Conversely $M$ itself
satisfies (E).

If an equivalence respects seeds and is stable under every native
program, equivalent states have equal membership in every generated
column, so it refines the constructed row equivalence. The latter is
itself stable by the first theorem. Neither part asserts that the
coarsest row quotient already supplies the least concept repair.
\end{proof}

\section{Three distinct tests for an unchanged characterization}\label{sec:repair}
For clarity take a finite carrier $X$, old columns $F\subseteq\Pow(X)$,
new seed columns $N$, and the \emph{entire resulting} native generator
set $A$, including old operations acting on newly introduced tests.
Old columns can already have been generated by the old operation set,
but that does not make old operations redundant on $N$. Write
\[
 F^+=\Gen_A(F\cup N),\quad M=\Meet(F),\quad M^+=\Meet(F^+).
\]
The object nerve $K_F$ has nonempty faces the nonempty subsets of an
actual column of $F$. Include the empty simplex separately.

\begin{proposition}[Exact repair tests]
\begin{enumerate}
\item $Q_{F^+}\to Q_F$ is bijective iff every new seed is constant on
old row classes and every $a\in A$ respects old row equality.
\item $M=M^+$ iff $N\subseteq M$ and
$a^{-1}(B)\in M$ for all $a\in A,B\in M$. It suffices to check
$B\in F$.
\item $K_F=K_{F^+}$ iff every nonempty new seed and every nonempty
$a^{-1}(C)$, for $a\in A,C\in F$, is contained in an old column.
\end{enumerate}
The second condition implies the first. Neither forces the third.
The third need not imply either of the others.
\end{proposition}
\begin{proof}
For (1), old-class constancy and operation stability propagate through
all words; conversely an unchanged row quotient is stable for all
new operations and observes all new seeds. For (2), inverse image
preserves intersections, so the stated generator check gives closure
of $M$ under all $A$-words. Minimality then gives equality. Necessity
follows from the stability of $M^+$.

For (3), a nonempty column is itself a face because $X$ is finite.
Thus unchanged complexes force each indicated column to be covered.
Conversely let $D(F)$ be the family of subsets covered by old columns,
together with the empty set. The one-step covering condition implies
$a^{-1}(S)\in D(F)$ whenever $S\in D(F)$, by choosing a column
$C\supseteq S$ and using
$a^{-1}(S)\subseteq a^{-1}(C)$. The extra seeds are in $D(F)$;
induction gives $F^+\subseteq D(F)$ and therefore equal nerves.

Every old extent is a union of old row classes, proving the implication.
The CNOT example below separates (1) from (2), and its reset stage
separates (2) from (3). Conversely, an old universal column already
makes the nerve a simplex; adding arbitrary separating columns can
change both rows and the lattice without changing that nerve.
\end{proof}
For infinite carriers, the nerve criterion must say that every finite
subset of each added column is old-covered. A single whole-column
cover need not follow. The exactness assertion here concerns equality
of simplicial complexes, not an equivalence of their homotopy types.



\section{Supported topology retains its own direction and domains}
\label{sec:supported}
For a finite context with column family $F$, let $K_F$ have the nonempty
faces $U$ contained in some actual column of $F$, and include the empty
simplex. Define the supported concept poset
\[
 \mathcal C_F=\{B\in\Meet(F): B\ne\varnothing,
                         \ \exists C\in F\ (B\subseteq C)\}.
\]
This is not obtained by deleting lattice extrema indiscriminately. The
nonempty face poset $\mathcal F(K_F)$ retracts onto $\mathcal C_F$ by
$U\mapsto\cl_F(U)$: the inclusion $s$ and closure $r$ satisfy
$rs=1$ and $1\leq sr$. Their order complexes are therefore homotopy
equivalent. This is the established supported FCA--Dowker bridge
\cite{Common,GP,BS}.

\begin{proposition}[Growth and partial operational transport]
For generated doctrines $G\subseteq G^+$, the forward comparisons are
$K_G\subseteq K_{G^+}$ and $\mathcal C_G\subseteq\mathcal C_{G^+}$.
On old faces the closures satisfy $\cl_{G^+}(U)\subseteq\cl_G(U)$;
this supplies an order homotopy, not an equality of the two closure routes.
For a shared program $p:i\to j$, its maximal supported domain is
\[
 K_p=\{\varnothing\}\cup
 \{\varnothing\ne U\subseteq X_i:X(p)[U]\in K_{G_j}\}.
\]
It is a subcomplex of $K_{G_i}$, and growth enlarges this domain.
For composable programs $p,q$ one has exactly
\[
 K_{qp}=\{U\in K_p:X(p)[U]\in K_q\}.
\]
The corresponding concept domain consists of nonempty $B\in M_i$ for
which $p_!(B)$ is supported. Closure intertwines these action spans,
including their domains. Backward projection $\rho^M$ is supported
only on new supported extents having an old common-test witness.
\end{proposition}
\begin{proof}
Old columns remain available and old extents remain closed, proving the
forward inclusions. New closure is smaller; both sides are nonempty and
have an old common witness, so they lie in the new supported poset.
An inequality of monotone maps induces the asserted order homotopy.
If $X(p)[U]\subseteq C\in G_j$, then
$U\subseteq X(p)^{-1}C\in G_i$, proving domain inclusion and also
showing that closing $U$ does not leave its domain. Equation (CN) gives
strict commutation with the concept action. If $X(qp)[U]$ has a
common column, pulling that column back along $q$ shows that $X(p)[U]$
already has a common column. This proves the domain composition formula;
the converse follows immediately. The concept version uses the same
argument and $q_!p_!=(qp)_!$. Finally $\cl_G(B)$ is old-supported
if and only if $B$ is contained in an old column. This is the exact
domain of the backward supported comparison.
\end{proof}
Growth yields commuting simplicial action spans and lax concept-action
inclusions. It need not yield partial maps with identical domains: an old
face may acquire a newly supported image. Nor does a partial action give
a global homology endomorphism. The rich-to-poor quotient and lattice
projection must not be confused with the poor-to-rich nerve inclusion.

\paragraph{A different mixed square can be strict.}
Evidence restriction is also distinct from language growth. If $E\subseteq F$
are columns on $H$ and $U\subseteq H$ is retained evidence, then for an
$F$-closed extent $A$,
\[
 \cl_E(A\cap U)\cap U\ \subseteq\ \cl_E(A)\cap U.
\]
Take $H=\{a,b,c\}$, $U=\{a,b\}$, old columns $\{a\},H$, and new
column $A=\{a,c\}$. The two sides are $\{a\}$ and $\{a,b\}$.
Both are supported and nonempty. If $U$ contains a representative of
every $F$-row, equality holds for every $F$-closed $A$: each member of
$A$ has a row-equivalent representative in $A\cap U$, so the two sets
have the same old common tests. This mixed right/left-adjoint comparison
is separate from the strict backward naturality proved above.


\section{When generated completion commutes with actual wiring}\label{sec:wiring}
Let $P\subseteq X\times Y$ be the \emph{actual} compatible joint-state
carrier, with projections $\pi_X,\pi_Y$. Component seed columns are
$E_X,E_Y$, and their native completions are $G_X,G_Y$.
Let $T$ be specified joint native generators $t:P\to P$; these may
be interacting operations, not products. Wire only the seeds first:
\[
 F_0=\{\pi_X^{-1}(C):C\in E_X\}
       \cup\{\pi_Y^{-1}(D):D\in E_Y\}.
\]
Completing separately before wiring gives
\[
 F_{\mathrm{loc}}=\{\pi_X^{-1}(C):C\in G_X\}
       \cup\{\pi_Y^{-1}(D):D\in G_Y\}.
\]
Completing the joint system gives $F_{\mathrm{joint}}=\Gen_T(F_0)$.
The relation between these families must be proved from native
capabilities, not inferred from a prospective product of quotients.

\begin{definition}[Component experiment liftability]
For each native component program $u$ on $X$ there is an admitted
joint word $\widetilde u$ with
$\pi_X\widetilde u=u\pi_X$; likewise for $Y$.
The equality concerns actual source transformations on all of $P$.
The lifting words and their resource admissibility are supplied native
certificates. They need not leave the other component unchanged.
\end{definition}
\begin{theorem}[Exact operational completion--wiring comparison]
Assume component experiment liftability. Then
\[
 F_{\mathrm{loc}}\subseteq F_{\mathrm{joint}}
 =\Gen_T(F_{\mathrm{loc}}).
\]
Consequently joint completion has the forced comparison directions
of the growth theorem. In particular put
\[
 L=\im\bigl(P\xrightarrow{(q_X\pi_X,q_Y\pi_Y)}Q_X\times Q_Y\bigr).
\]
The local row quotient is $L$, not an assumed entire product. Then:
\begin{enumerate}
\item The map $Q_{\mathrm{joint}}\twoheadrightarrow L$ is an
isomorphism iff every joint primitive respects the fibers of the
displayed profile map.
\item The local extent family is
\[
 M_{\mathrm{loc}}=
 \{P\cap(A\times B):A\in M_X,\ B\in M_Y\}.
\]
Full FCA completion commutes with wiring iff every joint primitive
pulls every such closed rectangle back to another such rectangle.
It suffices to check the two kinds of generating local columns.
\item The local and joint object nerves are equal iff every nonempty
one-step inverse image $t^{-1}(C)$ of a local column is contained in
a local column. Here the carriers are finite.
\end{enumerate}
If every joint primitive is componentwise native,
$t(x,y)=(u(x),v(y))$ on $P$, then the literal column families themselves
are equal. Combined with liftability this supplies a sufficient
source-level commutation theorem for all the characterizations.
\end{theorem}
\begin{proof}
For a generated component experiment $u^{-1}(C)$, its local cylinder
is $\widetilde u^{-1}(\pi_X^{-1}(C))$, hence belongs to the joint
completion. This proves the first inclusion. Since
$F_0\subseteq F_{\mathrm{loc}}\subseteq\Gen_T(F_0)$, applying the
idempotent generation operator gives the asserted equality.

Two joint states agree on all local tests exactly when their two
component quotient profiles coincide. This proves the description of
$L$ and (1) by the repair criterion. Intersecting any family of local
cylinders yields $P\cap(A\times B)$, and every such rectangle is an
intersection of local cylinders. This proves the displayed formula;
(2) and (3) now follow from the respective repair criteria. Finally a
componentwise native primitive pulls an $X$-column back along $u$ and
a $Y$-column back along $v$. Both are again local generated columns.
Thus $F_{\mathrm{loc}}$ is already $T$-stable, proving literal equality.
\end{proof}
The liftability assumption is substantive with shared stock, ownership,
feedback, or compatibility constraints. If it fails, even the displayed
inclusion of local experiments in joint experiments is not automatic.
Separate component characterizations may posit experiments that cannot
be executed together on the actual compatible carrier. That is a native
source obstruction, not a failure repaired by declaring a quotient map.

\section{A two bit interaction separates the three changes}\label{sec:cnot}
Take $X=Y=\{0,1\}$ and $P=X\times Y$, ordered $00,01,10,11$.
The component observations test $0$ and $1$; their initial native
operation categories contain only identities. The four local columns
are
\[
 x=0,\quad x=1,\quad y=0,\quad y=1.
\]
There are four distinct rows. The full extent family has ten sets:
$\varnothing$, four singleton sets, these four two-element columns,
and $P$. The object nerve is a square.

Now admit the actual reversible native operation
\[
 T(x,y)=(x,x\oplus y).
\]
Since $T^2=1$, all generated columns are exactly the four local columns
and
\[
 x\oplus y=0,\qquad x\oplus y=1.
\]
No additional column is concealed in a longer word. These six columns
are all six two-element subsets of $P$. The row quotient still has
four states, but the full extent family now has twelve sets: empty,
all four singletons, all six pairs, and $P$. The nerve is the graph
$K_4$, with no triangular faces, so
\[
 \beta_1(K_{\mathrm{local}})=1,
 \qquad \beta_1(K_{\mathrm{new}})=6-4+1=3.
\]
These are Betti numbers over $\mathbb F_2$, and also over every field.

The FCA obstruction is explicit. For $U=\{00,11\}$,
\[
 \cl_{\mathrm{local}}(U)=P,
 \qquad
 T[U]=\{00,10\},
 \qquad
 \cl_{\mathrm{local}}(T[U])=\{00,10\}.
\]
Thus $T[\cl_{\mathrm{local}}(U)]\not\subseteq
\cl_{\mathrm{local}}(T[U])$. No complete-join action on the old
concept lattice can agree with CNOT on its four marked object states.
The old state quotient is already discrete, so CNOT acts on it without
any repair. The twelve-element lattice is the least operational FCA
repair. This is a source-derived separation between the state and
concept versions of completion--wiring commutation.

Next admit the native reset $R(x,y)=(x,0)$. Closing under all mixed
$T,R$ words gives exactly
\[
 \{x=0,x=1,y=0,y=1,x\oplus y=0,x\oplus y=1,\varnothing,P\}.
\]
Indeed $R^{-1}(y=0)=P$, $R^{-1}(y=1)=\varnothing$;
$R^{-1}(x\oplus y=b)=\{x=b\}$;
$R^{-1}(x=b)=\{x=b\}$; and $T$ permutes the six pair-columns
and fixes $\varnothing,P$. Hence this eight-column family is closed
under both primitive inverse images and contains all one-step new
columns, proving exact all-word closure. The row quotient and extent
family remain four and twelve respectively: $\varnothing$ and $P$
were already extents. But the actual universal test $P$ makes the
nerve a full $3$-simplex, so $\beta_1$ drops from $3$ to $0$.

The prior common-generation manuscript already showed that a universal
column can change support topology without changing an unmarked full
lattice. The added content here is not that elementary phenomenon.
It is the independently specified native interaction, its complete
mixed-word experiment generation, the minimal FCA repair, and the
exact source-level separation between the three commutation criteria.



\section{Changing the formation of composite experiments}\label{sec:decomposition}
Fix a finite arity $n\geq1$. All indexing categories are small. Supply
categories $A_0,\ldots,A_n$, $B_0,\ldots,B_n$, functors
\[
 W:\prod_{i=1}^n A_i\longrightarrow A_0,\qquad
 V:\prod_{i=1}^n B_i\longrightarrow B_0,
 \qquad u_i:A_i\longrightarrow B_i\quad(0\leq i\leq n),
\]
and a natural transformation
\begin{equation}\label{eq:cell}
 \chi:V\Bigl(\prod_i u_i\Bigr)\Longrightarrow u_0W.
\end{equation}
The direction of this cell is essential. No inverse is assumed.
An arrow $V(u_i a_i)\to u_0W(a_i)$ is an admitted directed comparison,
not a license to execute its reverse.

A partial-characterization diagram here is a covariant functor
$X_i:B_i\to\Set$. Its elements may be labelled admitted experiments or
partial evidence, and arrows transport those elements along admitted
interface refinements. This variance is a premise. A contravariant
experiment semantics must be converted explicitly by using opposite
categories; it is not silently substituted. A model-vocabulary reduct
$R:\mathcal M_{\rm new}\to\mathcal M_{\rm old}$ has its own variance,
opposite to an old-to-new program compiler. A model-indexed family of
these diagrams can be reindexed along $R$ and the results applied pointwise.
No inverse to a model reduct or noninvertible observation decoder is
asserted.

Define composition by
\[
 \operatorname{Comp}_V(X_1,\ldots,X_n)
   =\Lan_V(\mathop{\boxtimes}_i X_i),\qquad
 (\mathop{\boxtimes}_iX_i)(b_1,\ldots,b_n)=\prod_iX_i(b_i).
\]
The pointwise formula uses the category $(V\downarrow b)$, whose objects
are $(b_i,f:V(b_i)\to b)$ and whose arrows $\beta$ satisfy
$f=f'V(\beta_i)$. Thus
\[
 \operatorname{Comp}_V(X_i)(b)
   =\mathop{\operatorname{colim}}_{(b_i,f)\in(V\downarrow b)}
            \prod_iX_i(b_i).
\]
Composition and its admitted decompositions are source data. This formula
does not prove physical independence, resource copying, or implementability.

\section{The directed comparison and its exact factorization criterion}
For $a\in A_0$, define
\begin{align*}
 K_a:(W\downarrow a)&\longrightarrow(V\downarrow u_0a),\\
 (a_i,k)&\longmapsto
       (u_i a_i,u_0(k)\chi_{(a_i)}).
\end{align*}
On arrows it applies each $u_i$. Naturality of $\chi$ verifies the comma
condition. The source diagram below is exactly the target diagram
precomposed by $K_a$. Consequently there is a canonical mate
\begin{equation}\label{eq:mate}
 \Theta_X:\Lan_W\bigl(\mathop{\boxtimes}_i u_i^*X_i\bigr)
       \Longrightarrow u_0^*\Lan_V\bigl(\mathop{\boxtimes}_i X_i\bigr).
\end{equation}
This is the direction ``change then compose'' to ``compose then change.''
It exists without exactness assumptions.

For $f:V(b_i)\to u_0a$, define $\Fact_a(b_i,f)$ to have objects
\[
 (a_i,h_i:b_i\to u_i a_i,k:W(a_i)\to a),
 \qquad f=u_0(k)\chi_{(a_i)}V(h_i).
\]
A morphism to $(a'_i,h'_i,k')$ is a tuple $\alpha_i:a_i\to a'_i$ with
\[
 h'_i=u_i(\alpha_i)h_i,\qquad k=k'W(\alpha_i).
\]
This category is exactly $((b_i,f)\downarrow K_a)$. Its connectedness
means connectedness of the underlying undirected graph; it does not make
its arrows invertible.

\begin{theorem}[Exactness detected by product representables]\label{thm:set}
The following are equivalent:
\begin{enumerate}[label=(\roman*)]
\item $\Theta_X$ is an isomorphism for every family $X_i:B_i\to\Set$.
\item Every $K_a$ is final.
\item Every $\Fact_a(b_i,f)$ is nonempty and connected.
\end{enumerate}
For the representable family $X_i=B_i(b_i,-)$, the fiber of
$\Theta_X(a)$ over $f$ is exactly $\pi_0\Fact_a(b_i,f)$.
\end{theorem}
\begin{proof}
The equivalence of (ii) and (iii) is the comma-category criterion for
finality. A final functor preserves colimits of every Set-valued diagram,
so (ii) implies (i).

Conversely take $X_i=B_i(b_i,-)$. Co-Yoneda gives
\[
 \Lan_V(\mathop{\boxtimes}_i B_i(b_i,-))(u_0a)
       \cong B_0(V(b_i),u_0a).
\]
Explicitly, the colimit representative
$(b'_i,g,x_i:b_i\to b'_i)$ maps to $gV(x_i)$. Every arrow $f$ is
represented by $(b_i,f,\id)$, and the arrows $x_i$ identify each other
representative with this canonical one. This establishes the claimed
bijection with the stated variance.

A source representative is $(a_i,k,h_i:b_i\to u_i a_i)$ and its image
is $u_0(k)\chi V(h_i)$. The generating identifications in the source
colimit are precisely the morphisms displayed in $\Fact$. They preserve
$f$, so the fiber over $f$ is the set of connected components of that
category. A bijective mate makes each such set a singleton, equivalently
nonempty and connected. This proves (iii) and the fiber assertion.
\end{proof}

An empty lift category gives a missing composite experiment. Multiple
components give distinct realizations of the same changed experiment
that no admitted refinement connects. Surjectivity alone does not imply
exactness. Theorem~\ref{thm:set} is standard exact-square mathematics:
compare the Yoneda and factorization criteria of
Batanin--Kock--Weber, Lemmas 2.2 and 2.4, and their Day-convolution
application, Theorem 2.11 \cite{BKW}. The application below, rather than
a renaming of cofinality, supplies the concrete experimental content.

\section{Truth, elements, and noninvertible witnesses}
\begin{proposition}[Boolean existence]
For $\Bool$-valued diagrams, composition by joins of meets is universally
exact if and only if, for every $a,b_i$, nonemptiness of
$B_0(V(b_i),u_0a)$ implies nonemptiness of
\[
 \coprod_{f\in B_0(V(b_i),u_0a)}\operatorname{Ob}\Fact_a(b_i,f).
\]
In thin categories this is exactly nonemptiness of each relevant
factorization category.
\end{proposition}
\begin{proof}
For necessity use the truth predicates that are true at $b'_i$ exactly
when $B_i(b_i,b'_i)$ is nonempty. A true target composite is equivalent
to existence of an arrow $V(b_i)\to u_0a$; a true source composite
is exactly the displayed existence of a lift. For sufficiency, a true
target value supplies a tuple of true input values at some $b_i$ and an
arrow $f$. The hypothesis supplies a lift of some arrow with these
endpoints. Monotonicity along its $h_i$ transports the input truths and
produces a true source composite. The mate gives the reverse implication.
\end{proof}

The qualification ``some arrow'' matters outside thin categories. Boolean
existence does not distinguish parallel $f$'s. Requiring a lift of every
individual $f$ would be too strong for this proposition.

For categorical witnesses, supply strict functors $E_i:B_i\to\Cat$ and
retain the Grothendieck category
\[
 \operatorname{WComp}_V(E_i)(b)
   =\int_{(b_i,f)\in(V\downarrow b)}\prod_iE_i(b_i).
\]
Its objects retain a decomposition and actual evidence objects; a morphism
retains a decomposition arrow $\beta$ and vertical witness arrows
$E_i(\beta_i)e_i\to e'_i$. These arrows may be noninvertible. The
functor $K_a$ induces a comparison $J_E(a)$ between the source and target
Grothendieck categories.

\begin{theorem}[Universal retained-witness exactness]\label{thm:witness}
For fixed $a$, $J_E(a)$ is an equivalence for every family
$E_i:B_i\to\Cat$ if and only if $K_a$ is an equivalence.
\end{theorem}
\begin{proof}
With all $E_i$ terminal, the total-category comparison is $K_a$ itself,
proving necessity. Conversely suppose $K_a$ is an equivalence. Full
faithfulness of the induced total functor follows from full faithfulness
on decomposition arrows and the unchanged vertical witness hom-sets.
For any target evidence object over $d$, choose an isomorphism
$K_a(c)\to d$. Transporting the evidence along the inverse isomorphism
provides a source object whose image is isomorphic to the given one.
Thus the total functor is also essentially surjective.
\end{proof}

Finality, or even contractibility of the factorization categories, is
not sufficient for Theorem~\ref{thm:witness}. The Grothendieck construction
is a specified witness-retaining semantics, not automatically the
Cat-valued left Kan extension and not a general physical-execution model.
No groupoid completion or inversion of witness arrows has been used.

\section{A complete classification for integer resource changes}\label{sec:resources}
Let $m,d\geq1$ and $n\geq2$. The poset category $\N^m$ has an arrow
$r\to s$ when $r\leq s$ coordinatewise and tensor given by addition.
A matrix $C\in\N^{d\times m}$ defines a strict symmetric monoidal
functor $r\mapsto Cr$. Resource units are indivisible. For covariant
diagrams $X_i:\N^d\to\Set$, set
\[
 (\mathop{\star}_iX_i)(s)
   =\mathop{\operatorname{colim}}_{s_1+\cdots+s_n\leq s}
      \prod_iX_i(s_i),
\]
with componentwise-order arrows among decompositions. The comparison is
\begin{equation}\label{eq:budget}
 \mathop{\star}_i(X_i\circ C)(r)
      \longrightarrow(\mathop{\star}_iX_i)(Cr).
\end{equation}
Use joins and meets for Boolean diagrams, or the preceding Grothendieck
categories for retained witnesses.

\begin{theorem}[Three exact resource criteria]\label{thm:matrix}
The following classify universal exactness of \eqref{eq:budget}, over
all budgets and all input diagrams of the stated doctrine.
\begin{enumerate}[label=(\roman*)]
\item Boolean exactness holds if and only if every column of $C$ is zero
or a standard basis vector.
\item Set-valued exactness holds if and only if, additionally, no two
nonzero columns are equal. Thus $C$ is a partial coordinate injection,
with zero columns allowed.
\item Universal equivalence of retained witness categories holds if and
only if, additionally, no column is zero. Thus $C$ is a coordinate injection.
\end{enumerate}
\end{theorem}
\begin{proof}
For a physical demand tuple $b_i$ with $\sum_i b_i\leq Cr$, the
factorization poset is
\begin{equation}\label{eq:alloc}
 \{(a_i):\ \sum_i a_i\leq r,\quad b_i\leq Ca_i\text{ for every }i\},
\end{equation}
ordered componentwise.

For (i), suppose a column $c=Ce_j$ has coordinate sum at least two.
Write $c=p+q$ with $p,q$ nonzero. At budget $r=e_j$, choose
$b_1=p,b_2=q$, and other demands zero. In any source allocation at most
one $a_i$ is nonzero, because $\sum_i a_i\leq e_j$. Both $b_1,b_2$
are nonzero, however, so both require nonzero allocations. Thus
\eqref{eq:alloc} is empty. A nonnegative integer column of coordinate
sum at most one is exactly zero or a basis vector.

Conversely suppose every column is zero or a basis vector. For each
output row $\ell$, let $J_\ell=\{j:Ce_j=e_\ell\}$. The demand
inequality says
\[
 \sum_i b_{i,\ell}\leq\sum_{j\in J_\ell}r_j.
\]
Distribute these integer supplies among the $n$ bins with prescribed
integer demands $b_{i,\ell}$, for instance greedily. This gives
$a_{i,j}\geq0$ with $\sum_i a_{i,j}\leq r_j$ and
$\sum_{j\in J_\ell}a_{i,j}=b_{i,\ell}$. Put zero in the zero-column
coordinates. This is an object of \eqref{eq:alloc}; the Boolean criterion
therefore proves (i).

For (ii), the empty-lift argument still applies. If two distinct columns
$j,k$ are $e_\ell$, choose $r=e_j+e_k$ and
$b_1=b_2=e_\ell$, with other demands zero. The lifting poset has exactly
two objects: give $e_j$ to the first component and $e_k$ to the second,
or the reverse. All supply is used; the objects are incomparable and
there are no other objects to connect them. For representables
$X_1=X_2=\N^d(e_\ell,-)$ and $X_i=\N^d(0,-)$ for $i>2$, the
comparison has two source elements and one target element.

Conversely, assume $C$ is a partial coordinate injection. A supported
row $\ell$ has a unique corresponding source coordinate $j$. Set
$a^{\min}_{i,j}=b_{i,\ell}$ and set every zero-column coordinate to zero.
Unsupported demand coordinates vanish because $\sum_i b_i\leq Cr$.
The tuple $a_i^{\min}$ is feasible and is below every other tuple in
\eqref{eq:alloc}. Hence the lifting poset has a least object and is
nonempty and connected. Theorem~\ref{thm:set} proves (ii).

For (iii), universal witness equivalence implies that every decomposition
functor is an equivalence, hence final; thus (ii) is necessary. If column
$j$ is zero, take $r=e_j$. The target decomposition category at $Cr=0$
is terminal. The source has the all-zero allocation and the distinct
allocation that gives $e_j$ to component one, with a noninvertible arrow
between them. This poset has nonisomorphic distinct objects and is not
equivalent to a terminal category. Hence zero columns are forbidden.

Conversely, if the columns are distinct basis vectors, each target
decomposition below $Cr$ has exactly one inverse source decomposition,
obtained by reading its supported coordinates. Unsupported target
coordinates are forced zero. The inverse preserves and reflects order,
so the decomposition functor is a poset isomorphism. The induced total
witness-category comparison is an isomorphism, up to product conventions.
\end{proof}

\subsection{Three small experiments and a positive case}
\paragraph{Batch splitting.}
For $C=(2)$ at $r=1$, two one-unit experiments fit a total physical budget
of two. Each changed local interface, however, uses a whole source packet;
the local packet allocations cannot serve both components. The lifting
category for $(1,1)$ is empty. A sealed packet of two calibration reagents
is a physical interpretation only if the source local interface actually
forbids inter-component splitting and the target global interface permits
it. Those access conditions, not the numeral two alone, create the example.

\paragraph{Forgotten provenance.}
For $C=(1\ \ 1)$ and $r=(1,1)$, two labelled stocks become one physical
reagent type. Two one-unit experiments can always be supplied, but their
two alternative stock assignments form a discrete two-object lifting
category. Boolean feasibility agrees; element-level allocation data does not.

\paragraph{Invisible audit resources.}
For $C=(1\ \ 0)$ at $r=(0,1)$, the target decomposition category is
terminal. In the source audit coordinate it is the span
\[
 (0,0)\longrightarrow(1,0),\qquad
 (0,0)\longrightarrow(0,1).
\]
The constant-diagram colimit is one point and the functor to the terminal
category is final. It is not an equivalence: the displayed arrows are
noninvertible and the three objects are not isomorphic. In fact this
$C$ is exact for every Set-valued diagram, not just for this constant one,
by Theorem~\ref{thm:matrix}. Thus the witness obstruction survives even
under universal extensional exactness.

\paragraph{A positive vocabulary change.}
The map $C(q)=(q,0)$ changes the resource vocabulary by introducing an
additional target resource unavailable in image budgets. Every decomposition
of an image budget has zero in that coordinate and lifts uniquely. It is
exact even at the retained-witness level. The map $C(q,a)=(q,0)$ is a
nontrivial Set-exact change that additionally forgets an audit coordinate,
but it is not witness-exact.

\subsection{Not the earlier manufacture theorem}
The programme's adaptive-manufacture theorem compares
$\bigvee_h Cd_h$ with $C(\bigvee_h d_h)$; its universal criterion is
that each row has at most one nonzero entry. That permits arbitrary
positive weights and permits one source column to use multiple physical
resources. In contrast, $C=(2)$ satisfies that earlier criterion but
fails even Boolean convolution exactness. The matrix $C=(1\ \ 1)$
fails the earlier criterion but satisfies Boolean convolution exactness.
The new calculation concerns decomposing a fixed budget, whereas the old
one concerns joining alternative future demands. They are different
interchanges, not alternative terminology for the same theorem.

\section{The exact response bridge to heterogeneous characterizations}\label{sec:slice}
Let $S$ be a supplied actual system-state set. Write $\mathcal P(S)$
for the set of response columns, regarded as a discrete constant diagram.
Suppose the target decomposition diagram $D$ at an output interface has
a natural cocone
\[
 \epsilon_d:D(d)\longrightarrow\mathcal P(S),
\]
and equip the source diagram $DK_a$ with precisely the pulled-back cocone.
Naturality means that an admitted refinement preserves the response
column of each transported label. This is an extra semantic requirement;
a stateful intervention that changes the experiment's meaning does not
automatically satisfy it.

\begin{proposition}[Response-preserving comparison]\label{prop:slice}
The induced colimit comparison is a morphism in $\Set/\mathcal P(S)$.
If it is bijective, it gives an isomorphism of labelled incidence contexts,
with identity map on $S$. Consequently it preserves the induced behavioral
state equivalence, the formal-concept lattice, and the supported Dowker
complexes up to their canonical isomorphisms. For the same supplied deterministic state and program interpretation, if the
labelled tests carry specified program-precomposition actions and the
comparison is equivariant, the corresponding quotient and concept actions,
and the object-nerve action spans with their stated domains, transport as well.

More weakly, equality of the sets of evaluated response columns is enough
for the behavioral state equivalence, extent-based concept lattice,
supported concept core, and object-side Dowker complex to agree. In
particular, surjectivity of the response-preserving comparison suffices
for these extensional conclusions; injectivity is not needed.
\end{proposition}
\begin{proof}
The universal property of each colimit gives evaluation maps to
$\mathcal P(S)$, and the pulled-back-cocone hypothesis makes their triangle
with $\Theta$ commute. A bijection over $\mathcal P(S)$ preserves exactly
the incidence relation $s\mathrel I t$ defined by $s\in\epsilon(t)$.
All the stated constructions are invariant under an incidence isomorphism:
rows have the same equalities, derivation operators correspond under the
test bijection, and common-test simplices correspond under the same
bijection and identity on states. Equivariance identifies the induced quotient and concept actions by their
defined formulas. The common supplied state maps give the same vertex maps
on the object nerves, and the common response columns give their domains.
This does not infer equality of raw state maps from nonseparating tests.
For the weaker conclusion, duplicate labels do not change row equality,
intersections of response columns, whether a closed extent has a
nonempty common-test set, or whether a finite state set has a common
test. Surjectivity over $\mathcal P(S)$ guarantees equality of the
sets of response columns. It does not give a bijection of labelled tests
or an isomorphism of their witness categories.
\end{proof}

For independent parallel tests, a concrete evaluation is
$\pi_X^{-1}C_e\cap\pi_Y^{-1}C_f$ on the admitted joint state set.
That example still requires actual independent execution and refinements
preserving this response. A nonbijective label comparison need not alter
any state quotient or concept lattice: distinct labels can have equal
response columns. Thus Theorem~\ref{thm:matrix} does not claim that every
resource mismatch changes each downstream invariant.

The companion interaction-closure result uses tests pulled back along an
actual interacting operation, such as a reversible gate. Its native test
syntax is contravariant under program precomposition, not automatically a
Day-convolution diagram. Proposition~\ref{prop:slice} is the precise
conditional bridge; no claim is made that every interacting gate already
has the Kan-extension composition used here.


\paragraph{A response-level realization of the batch obstruction.}
In the batch example $C=(2)$ and source budget $1$, let each component
have a single experiment $e_i$, available at physical budget at least $1$.
Let the actual joint state set be $S=\{0,1\}^2$, and let $e_i$ ask
whether its component bit is $1$. At physical budget $2$ the joint
experiment has response column $\{(1,1)\}$, given by the intersection
of the two cylinders. After local interface change, each experiment needs
one whole source packet, so no joint label exists at source budget $1$.
Restrict attention to this prescribed paired-test family: no zero-budget
no-op test was supplied to either input diagram. Individually executable
probes on the joint carrier are not being counted. The source paired-test
family is empty; the target has one response column. Their row quotients have respectively one and two classes, their
extent families one and two elements, and their object nerves respectively
no vertices and one vertex. This is an actual incidence consequence of
an absent decomposition, under the stated independent-test interpretation.
By contrast, the two disconnected provenance allocations for $C=(1\ \ 1)$
can evaluate to the same response column. They then obstruct label-level
Set exactness without changing any of these extensional invariants.

\section{Coherence of successive interface changes}
For a second comparison $(B,V)\to(D,Z)$ with maps $v_i$ and
$\psi:Z(\prod_i v_i)\Rightarrow v_0V$, the composite cell is
\[
 Z(\prod_i v_i u_i)
  \xRightarrow{\psi_{(u_i-)}}v_0V(\prod_i u_i)
  \xRightarrow{v_0\chi}v_0u_0W.
\]
Its decomposition functor is exactly
$K^v_{u_0a}K^u_a$. The induced mates therefore compose after the evident
reindexings. On a representative, both routes have composite arrow
$v_0u_0(k)\,v_0\chi\,\psi$. Identities and three-stage associativity
follow from the same formula. Final functors compose, as do equivalences,
so exact interface changes compose in the corresponding doctrine.

This is a theorem about a supplied directed composition square. It does
not manufacture an operad, feedback semantics, contextad, physical
compiler, native quotient, or resource reset. If those richer structures
are intended, their units, substitutions, witness transport and coherence
must be supplied and checked. The prior fixed-interface regular-image
result starts from a fixed test semantics and compares quotients of state
carriers; the present result changes which joint experiments can be formed
before taking any state quotient. The two problems are complementary.



\section{Native synchronization exposes the missing resource correlation}
\label{sec:native}
The decomposition category in the preceding theorem retains formal
allocation witnesses. A separate operational interpretation is required
before such a witness is an executable run. We now give one concrete
native class, its sharp obstruction, and a primitive certificate that
compiles local witnesses into joint witnesses. Its composition is
synchronization, not the additive-budget Day convolution above.

Let $B,A$ be finite concrete and visible alphabets, and let
$\alpha:B\twoheadrightarrow A$ be a letterwise, nonerasing interface
map. A component is a finite directed multigraph of actual configurations,
with current state $x$, terminal set $F$, and edges
$e:s\xrightarrow{b}t$ carrying costs in $\N$. Edges are controllably
selectable, terminating native commands or certified single-event macros.
Endpoints retain residual resources and installed state. A run is a
composable edge sequence; its cost is the sum of edge costs.
Define its concrete trace profile
\[
 f_{M,x}(w)=\min\{c(p):p:x\leadsto F,\ \ell(p)=w\}\in\Cost,
 \qquad \min\varnothing=\infty.
\]
Finite values have native run witnesses, found by finite-word dynamic
programming. This is complete-table controllable planning, not sure
control against adversarial outcomes or empirical identification of an
unknown machine.

Independent parallel components synchronize the same concrete letter at
each step and add their costs. Independence here is an actual product
runtime with separately accounted resources; shared stock must instead
be represented once in a joint state. For profiles $f_1,\ldots,f_m$,
interface change gives
\[
 (\alpha_!f_i)(v)=\min_{\alpha(w)=v}f_i(w).
\]
Changing each component first and then adding its local optimum yields
$L(v)$, while changing an actual synchronized composite yields $J(v)$:
\[
 L(v)=\sum_i\min_{\alpha(w)=v}f_i(w),\qquad
 J(v)=\min_{\alpha(w)=v}\sum_i f_i(w).
\]
Their difference is the requirement of a common concrete word.
The abstract number $L(v)$ is not itself an executable joint witness.

\begin{theorem}[Exact optimal-witness obstruction]\label{thm:native-gap}
For every visible word $v$, $L(v)\leq J(v)$. If all local minima are
finite, put
\[
 A_i(v)=\{w:\alpha(w)=v,\ f_i(w)=(\alpha_!f_i)(v)\}.
\]
Then $J(v)=L(v)$ if and only if $\bigcap_iA_i(v)\ne\varnothing$.
A common minimizing word together with its native component runs gives
a synchronized witness of cost $L(v)$. If any local minimum is infinite,
both sides are infinite.
\end{theorem}
\begin{proof}
For each eligible $w$, every $f_i(w)$ is at least its own fiber minimum.
This proves the inequality. Since $B$ is finite and the map is letterwise,
the fiber over $v$ is finite. When equality is finite a minimum is attained;
the sum of the nonnegative differences from the local minima is zero,
so every difference is zero. Conversely a common minimizer attains $L(v)$.
The runs synchronize step by step because they have the same concrete
word and independent actual endpoints. If one fiber has no finite run,
no jointly finite run exists.
\end{proof}

The trace-and-budget predicate $(v,r)$ says that a terminal native run
of visible word $v$ exists with cost at most $r$. Treat these as semantic
feasibility predicates unless an admitted decision readout is supplied,
for example a public current configuration with complete certified tables
and permitted decision computation. An existential success has a native
run witness; failure of a search is not automatically a physical observation
on an unknown state. Under this declared test doctrine, evaluate the
predicates on actual state tuples to obtain response columns. Values in $\Cost$ are determined
by all their natural-number thresholds. Therefore the two routes have
exactly the same budget-test incidence if and only if $J=L$ on every
state tuple and word. A discrepancy $L(v)<J(v)$ is witnessed by the
single threshold $(v,L(v))$. Under equality, the response-column lemma
of Section~\ref{sec:slice} gives equality of the resulting row, concept
and object-nerve constructions. Under inequality, the incidence differs;
not every downstream invariant must detect that difference.

\begin{example}[Boolean exactness and unbounded budget error]
Let $B=\{0,1\}$, $A=\{a\}$, and hide both letters as $a$.
Two one-state cyclic components have both loops, with their sole state
terminal. The costs of $0,1$ are $1,K+1$ in the first component and
$K+1,1$ in the second, where $K\geq1$. Then for every $n\geq0$,
\[
 L(a^n)=2n,\qquad J(a^n)=(K+2)n.
\]
Every visible trace is feasible by both routes. Locally optimal witnesses
are $0^n$ and $1^n$; for $n>0$ they cannot synchronize. Every common
concrete word instead has total cost $(K+2)n$. Thus feasibility, terminal
conditions and even complete unbounded visible language agree while the
budget defect grows without bound. This is neither a missing-state
artifact nor a failed heuristic search for a witness.
\end{example}

\begin{proposition}[Sequencing is exact for the same interface change]
For weighted profiles, define
\[
 (f*g)(w)=\min_{w=uz}(f(u)+g(z)),\qquad
 (f\wedge g)(w)=\min(f(w),g(w)).
\]
Then $\alpha_!$ preserves $*$, $\wedge$, the sequential unit and
finite iteration $f^*=\min_{n\geq0}f^{*n}$.
These equalities retain native witnesses when concatenation uses the same
actual intermediate endpoint. For variable endpoints the exact formula is
\[
 (F*G)_{s,t}(w)=\min_{r,\,w=uz}
                      (F_{s,r}(u)+G_{r,t}(z));
\]
entrywise pushforward preserves it as well.
\end{proposition}
\begin{proof}
Both convolution routes minimize $f(u)+g(z)$ over concrete words with
$\alpha(u)\alpha(z)=v$. A concrete split gives a visible split,
and lifts of a visible split concatenate. The intermediate-state minimum
can be exchanged in the same way when that state is retained. Choice
exchanges minima; the unit is preserved because $\alpha$ is nonerasing.
Iterate convolution and exchange the minima over words and iteration
number. Every finite optimum has a finite native run witness.
\end{proof}
Thus the synchronization failure remains even though the same interface
change preserves choice, native sequencing and arbitrary finite repetition.

\begin{theorem}[Universal synchronization contexts]\label{thm:universal-native}
For $f:B^*\to\Cost$, the identity
\[
 \alpha_!(f+g)=\alpha_!f+\alpha_!g
\]
holds for every $g:B^*\to\Cost$ if and only if $f$ is constant on
each fiber of $\alpha$. It suffices to test the identity against the
finite native exact-word testers $\delta_z$, with value $0$ at $z$
and $\infty$ elsewhere.
\end{theorem}
\begin{proof}
A finite zero-cost chain spelling $z$, with only its final state terminal,
realizes $\delta_z$. At $v=\alpha(z)$ the stated identity reads
$f(z)=(\alpha_!f)(v)$. Requiring this for every $z$ says exactly that
$f$ is fiberwise constant, including infinite values. Conversely a
constant value $c$ on the fiber can be taken outside the minimum of
$c+g(w)$; if $c=\infty$ both sides are infinite.
\end{proof}
This is stronger than saturation of a Boolean trace language. It identifies
precisely what information must be retained for arbitrary concrete
synchronization contexts. A selected compatible class of components can
satisfy the weaker constructive condition that follows.

\begin{definition}[Common endpoint-preserving normalization]
Choose one section $\kappa:A\to B$ of $\alpha$, shared by every
component. For every native edge $e:s\xrightarrow{b}t$, supply an
effective selectable native edge or certified single-event macro
\[
 N(e):s\xrightarrow{\kappa\alpha(b)}t,
 \qquad c(N(e))\leq c(e).
\]
Source and target are the same actual configurations. Internal macro steps
introduce no extra synchronization events, and their termination, cost
and endpoint are separately certified. No hidden information is required
to choose the replacement.
\end{definition}

\begin{theorem}[A primitive native compiler for exact change and wiring]
\label{thm:native-normal}
Under common endpoint-preserving normalization, for every state tuple,
visible word and finite number of independent components,
\[
 (\alpha_!f_i)(v)=f_i(\kappa(v)),\qquad J(v)=L(v).
\]
Every family of locally chosen native witnesses with visible word $v$
compiles to synchronized native witnesses with concrete word $\kappa(v)$,
the same individual endpoints and no larger individual costs. The compiler
commutes with concatenation and works for arbitrary finite runs in cyclic
machines.
\end{theorem}
\begin{proof}
Replace every edge and concatenate. Identical actual endpoints make the
result a legal run from the same initial configuration to the same final
configuration. Its word is $\kappa\alpha(w)$ and its cost does not
increase. Hence $f_i(\kappa(v))\leq f_i(w)$ for every lift $w$ of $v$.
Since $\kappa(v)$ itself is such a lift, it realizes the fiber minimum.
All components now share a minimizing word, proving equality by
Theorem~\ref{thm:native-gap}. Occurrencewise replacement commutes with
concatenation. Normalized runs synchronize because their canonical letters
match at every step. None of these arguments bounds the number of loop
iterations or identifies merely observationally equivalent endpoints.
\end{proof}

For finite graphs without internal macros, this is checked by a finite
scan of actual edges. More sharply, put
$M_b(s,t)=\min\{c(e):e:s\xrightarrow b t\}$. With every actual state
allowed as a named terminal endpoint, canonical words realize every fiber
minimum if and only if
\[
 M_{\kappa\alpha(b)}(s,t)\leq M_b(s,t)\quad\text{for all }s,t,b.
\]
Likewise every endpoint-specific profile is fiberwise constant if and only
if $M_b=M_{b'}$ whenever $\alpha(b)=\alpha(b')$.
Necessity uses one-letter tests. Sufficiency follows by replacing each
witnessed edge, or by monotonicity and equality of the corresponding
min-plus matrix products. Thus both the restricted common-normalization
class and the universal-context class have independent finite primitive
certificates.

The two cyclic machines above have different cost-reducing canonical
letters, so their separate normalizers do not compose. Exact endpoint
preservation and a \emph{common} section are the substantive conditions.
Vector resource costs admit the same constructive theorem when every
replacement is componentwise no more expensive. A scalar minimum cannot
otherwise represent a Pareto frontier, and pooling total budget does not
by itself preserve each component's contractual allocation.


\section{The precise contextad connection}\label{sec:contextads}
Contextads organize the operational side of this comparison, provided their
actual structural hypotheses are supplied. They do not extract an
observation doctrine from a table. In the current source, Capucci--Myers
\cite[Definition 5.25, Examples 5.33--5.35]{CM2}, a contextad has indexed
grades $K:C^{\mathrm{op}}\to\Cat$, extension $e:\int K\to C$,
unit grades and coherent accumulation. A contextual arrow $(P,f):A\Rightarrow B$
has body $f:e_AP\to B$. The grade of a composite with $(R,g):B\Rightarrow D$
is $P\otimes f^*R$.

A strong contextad morphism supplies $F:C\to C'$, a cartesian functor
$G:\int K\to\int K'$ over $F$, and a natural lineator
\[
 \ell_{A,P}:e'_{FA}(G_AP)\longrightarrow F(e_AP),
\]
together with coherent invertible unit and multiplication comparisons
\[
 G_AI_A\cong I'_{FA},\qquad
 G_A(P\otimes Q)\cong
 G_AP\otimes'\ell_{A,P}^*G_{e_AP}Q.
\]
The cartesian condition means preservation of cartesian lifts and hence
$G(u^*R)\cong(Fu)^*G(R)$; it does not say that the entire square of
categories is a pullback. Compatibility with the extension structure,
units and associativity is required. The lineator itself need not be
invertible. The resulting strong operational comparison is
\[
 T(P,f)=(G_AP,F(f)\ell_{A,P}).
\]
Indeed the displayed grade comparison and cartesian property identify
$G_A(P\otimes f^*R)$ with
$G_AP\otimes'(F(f)\ell_{A,P})^*G_BR$; the lineator compatibility
identifies the bodies. Unit and associativity are the stated coherences.
This is the established construction of \cite[equations (5.6.20)--(5.6.21)]{CM2}.

For the characterization bridge, an additional interpreted state map
$f_a:X_a\to Y_{Fa}$ must intertwine each admitted source operation $u$
with its translated operation $\widetilde u$:
$f_bu=\widetilde u f_a$. There must also be a backward translation of
\emph{admitted} labelled tests $g_a:D_{Fa}\to E_a$, satisfying
\[
 f_a(x)\models d\ \Longleftrightarrow\ x\models g_a(d),
 \qquad g_a\widetilde v_u=v_ug_b,
\]
where $v_u,\widetilde v_u$ are the respective continuation operations
on labels. This is the exact Chu condition with operation coherence.
It requires every pulled-back target column to be represented by an
actual source test, not merely by a semantic intersection of source columns.
Coherent label choices are stronger than columnwise representability.

\begin{proposition}[Operational and observational hypotheses together]
Under these state and test equations, $[x]\mapsto[f_a(x)]$ is a
well-defined operation-equivariant row map. Inverse image of extents has
the canonical left adjoint
\[
 (f_a)_!(A)=\cl_D(f_a[A]),\qquad (f_a)^*(B)=f_a^{-1}B.
\]
It preserves marked object concepts and intertwines the lifted operations.
The simplicial comparison has exactly the domain
\[
 K_{f_a}=\{\varnothing\}\cup
 \{\varnothing\ne U\in K_E:\exists d\in D_{Fa},\ U\subseteq C_{g_a(d)}\}.
\]
These comparisons compose, with the composite domain equal to the domain
of the first map whose image lies in the domain of the second.
\end{proposition}
\begin{proof}
Source-row equality implies equality on every $g_a(d)$, hence on every
target test after $f_a$. Preimages of target extents are intersections
of the columns $C_{g_a(d)}$, so they are source extents. This proves the
adjunction and its marked-generator property as in (CN). State and test
coherence give operation equivariance. A finite image $f_a[U]$ has a
common target witness precisely when the displayed domain condition holds.
For a second test translation $h$, the composite condition is the
existence of a witness in $\im(gh)$; this is exactly membership in the
successive domains. Closing a face preserves its common witnesses,
so the corresponding supported concept comparison has the same domain.
\end{proof}

In the display-map instance, contextual composition is an actual pullback
of spans. A functor $F$ must preserve the required display pullbacks for
a strong comparison. Otherwise the canonical map is only
\[
 F(X\times_BY)\longrightarrow FX\times_{FB}FY.
\]
With forward apex maps as span cells, this is the \emph{oplax} direction.
Neither finality of an unrelated Set-valued decomposition functor nor a
row quotient supplies its inverse. This is the exact point at which a
contextad-based application needs a new source-level proof.

The cited v2 is dated 3 November 2026 and is titled \emph{Contextful
Arrows via Wreath Products}. It is not the earlier v1 with its earlier
title and numbering. Its Section 5.4 and Remark 5.20 retain qualifications
about the free Kleisli tricocompletion and higher wreath formula. The
concrete grade calculation above does not declare those qualifications
settled, and none of the main results in this note depends on doing so.


\section{What is inherited and what this note establishes}\label{sec:prior}
The distinction matters both mathematically and historically. A general
cofinality theorem, an ordinary partition refinement, or a renamed Chu
adjunction would not by itself advance the present question.

\paragraph{Exact squares.}
Theorem~\ref{thm:set} is the standard Guitart exactness criterion.
Batanin--Kock--Weber \cite[Lemmas 2.2 and 2.4]{BKW} prove the Yoneda
and connected-factorization tests; their Theorem 2.11 gives the related
Day-convolution and monoidal Kan-extension theorem under its stated
hypotheses. They attribute the underlying criterion to Guitart (1980).
That original paper was not independently checked here. The present
contribution at this level is the explicit integer-budget classification,
its three different failure mechanisms, and its response-semantic and
native-execution boundaries. No new general categorical exactness principle
is claimed, and historical priority for the matrix specialization has not
been established.

\paragraph{Concepts and topology.}
Ern\'e \cite[Theorem 7.1]{Erne} gives marked-context representation via
doubly based lattices; the markings are essential. Gratus--Porter
\cite[Proposition 1(ii), Corollary 1]{GP} already restrict Chu nerve maps
to transported attribute witnesses. Brun--Salbu \cite[Definition 3.2,
Theorem 4.3]{BS} give the natural rectangle-complex comparison for
covariant relation maps. Chowdhury--M\'emoli \cite[Theorem 3,
Corollary 20, Proposition 15]{DowkerPersistence} establish functorial
Dowker persistence and a metric stability theorem. Arbitrary test changes
are not automatically their filtrations or satisfy their metric hypotheses.
The audit verified these antecedents, not an exhaustive literature on
incremental concept-lattice algorithms.

\paragraph{Recognition and vocabulary change.}
Clarke et al. \cite[Lemma 4.12]{Optics} state enriched double Yoneda;
Theorem 4.14 applies it to generalized Tambara modules and optics.
Recognition by all copresheaf evaluations is therefore not a new principle.
Finite available probes or a restricted native grammar need their own
completeness theorem. Institutions \cite[Definition 1, Theorem 11]{Institutions}
already make sentences covariant, models contravariant and translated
theories compositional. Neither institution satisfaction nor a contextad
morphism alone proves physical probe availability.

\paragraph{Learning with changed interfaces.}
CALF \cite[Theorems 16 and 20]{CALF} gives hypothesis and testing
criteria under explicit closedness, consistency, reachability and
observability assumptions. Henry et al. \cite[Theorem 5.3]{CoalA}
construct compositional active learning with alphabet refinement under
teacher and local-learner assumptions. Their Remark 5.4 does not guarantee
recovery of the original component count or each original component.
Rosebrugh--Sabadini--Walters \cite[Propositions 43 and 46,
Corollary 47]{MinimalRealization} already organize minimization after
composition by an idempotent lax construction. The present finite-source
calculations do not supersede these learning or minimization results.

\paragraph{The three preceding project manuscripts.}
The common-generation note \cite{Common} already proves the fixed acted
incidence construction, the canonical marked concept action, the supported
FCA--Dowker comparison and its exact partial domains. It already includes
the universal-test obstruction. The wiring note \cite{Wiring} retains
the actual realizability image and separates joint lifting from exterior
separation. Its later three-level adequacy theorem separates row coverage,
intersection generation and support, with a separate result for operation
domains. The adaptive-experiment note \cite{Adaptive} already proves
manufacture-stable partition refinement and native resource-preserving
policy transfer from the same residual configurations. Its resource
interchange compares a join of alternative future demands with resource
translation, not a decomposition of one fixed budget.

Accordingly, the present note does not rebrand their fixed-interface
conclusions. It derives simultaneous changes of tests and operations,
compares their backward and forward transport, and computes a native
interaction where the state, concept and support criteria separate.
The general exact-square theorem is then applied to a different upstream
question: which joint experiments can be formed after a budget vocabulary
change? The matrix trichotomy answers that question exactly. The native
trace results refine the old existence obstruction by distinguishing
common feasible witnesses from common optimal witnesses and by providing
a primitive, endpoint-preserving compiler. These are proved mathematical
increments and a synthesis; no exhaustive novelty certification is asserted.

\section{Scope and the next concrete question}\label{sec:limits}
The supplied state semantics, admitted programs, initial observations,
resource vocabulary and wiring are part of the experimental contract.
The results do not infer those promises from a finite sample, invent an
unrestricted operational language, or identify every row by one destructive
experimental policy. Extensional conjunctions are not free resource-sharing
instructions. New observations may also require old operations to be
reanalyzed on new tests: simultaneous closure is under the \emph{entire}
resulting signature. A one-pass separate repair need not suffice.

The sharp hypotheses are deliberately at different levels. Rectangle
preimage stability concerns the concept semantics. Finality concerns
Set-valued colimits of supplied experiment diagrams. Equivalence of their
Grothendieck categories concerns retained formal witnesses. Native
normalization concerns actual run endpoints and costs. A proof at one
level must pass through a response-preserving evaluation or an explicit
compiler before it can be used at another.

A focused next problem is to weaken the common-section certificate in
Theorem~\ref{thm:native-normal} while retaining finite checkability and an
actual endpoint-preserving compiler. The desired class would allow the
choice of a common concrete letter to depend on a finite, publicly known
synchronization state, rather than one fixed $\kappa$. It should prove
optimal budget preservation for every finite run, retain the residual
state after each step, and identify when local certificates compose over
cycles. The cyclic counterexample supplies a nonnegotiable regression
case, and Theorem~\ref{thm:universal-native} identifies the stronger
arbitrary-partner boundary. This is a specific native synthesis question,
not a claim already discharged by categorical exactness.

\section{Verification and accompanying sources}\label{sec:verification}
All universal claims used in the note have proofs above. Independent finite
checks are regression evidence, not proofs of the unrestricted statements.
The accompanying scripts check:
\begin{itemize}
\item all distinct-column contexts on at most three objects with all unary
maps: 6,980 cases, 20,940 exactness comparisons and 173,796 refinement/action
squares; the all-word CNOT/reset closures, counts and homology;
\item 614 integer matrices in the stated bounded families, 5,696 budgets
and 146,466 factorization fibers for the Boolean/Set/witness trichotomy;
\item 85,219 scalar synchronization cases, 4,096 finite partner tests,
480 convolution values and cyclic trace families through length 20;
\item adversarial examples involving mixed evidence/language changes,
shared resources, cyclic joint feasibility and topology under logical closure.
\end{itemize}
A separate Lean 4.35.0/mathlib check proves only the elementary
row-quotient surjection for test inclusion, its naturality for a shared
stable operation, and an explicit strictness obstruction. Its reported
axioms are \texttt{Quot.sound} and \texttt{propext}; it has no
\texttt{sorry}. This does not mechanize the concept, topology,
Kan-extension, resource or native-normalization theorems.

The TeX source has been checked for balanced environments, label/citation
references and ordinary source integrity. It has intentionally not been
compiled. The companion archive preserves the component research drafts,
correction notes, source audits, scripts, check outputs, recovered preceding
manuscripts and assembly checkpoints. Internal access receipts and unrelated
files are not required to understand or reproduce the mathematics.


\begin{thebibliography}{99}
\bibitem{BKW} M. Batanin, J. Kock and M. Weber,
\emph{Regular patterns, substitudes, Feynman categories and operads},
Theory and Applications of Categories 33(7) (2018), 148--192.
\url{https://webercat.au/rpfc.pdf}.
\bibitem{Optics} B. Clarke, D. Elkins, J. Gibbons, F. Loregian,
B. Milewski, E. Pillmore and M. Rom\'an,
\emph{Profunctor Optics: a Categorical Update}, Compositionality 6(1) (2024).
\url{https://doi.org/10.32408/compositionality-6-1}.
\bibitem{Erne} M. Ern\'e, \emph{Categories of Contexts},
arXiv:1407.0512v1 (2014).
\url{https://arxiv.org/html/1407.0512v1}.
\bibitem{GP} J. Gratus and T. Porter, \emph{A geometry of information, I},
Dagstuhl Seminar Proceedings 04351 (2005).
\url{https://drops.dagstuhl.de/storage/16dagstuhl-seminar-proceedings/dsp-vol04351/DagSemProc.04351.6/DagSemProc.04351.6.pdf}.
\bibitem{BS} M. Brun and M. Salbu, \emph{The Rectangle Complex of a Relation},
Mediterranean Journal of Mathematics (2023).
\url{https://doi.org/10.1007/s00009-022-02213-0}.
Primary manuscript: \url{https://arxiv.org/pdf/2207.02018}.
\bibitem{DowkerPersistence} S. Chowdhury and F. M\'emoli,
\emph{A functorial Dowker theorem and persistent homology of asymmetric networks},
Journal of Applied and Computational Topology (2018).
\url{https://doi.org/10.1007/s41468-018-0020-6}.
\bibitem{Institutions} J. A. Goguen and R. M. Burstall,
\emph{Institutions: Abstract Model Theory for Specification and Programming},
Journal of the ACM (1992).
\url{https://doi.org/10.1145/147508.147524}.
\bibitem{CM2} M. Capucci and D. J. Myers,
\emph{Contextful Arrows via Wreath Products}, arXiv:2410.21889v2,
3 November 2026.
\url{https://arxiv.org/pdf/2410.21889v2}.
\bibitem{CALF} G. van Heerdt, M. Sammartino and A. Silva,
\emph{CALF: Categorical Automata Learning Framework}, CSL 2017.
\url{https://doi.org/10.4230/LIPIcs.CSL.2017.29}.
\bibitem{CoalA} M. Henry, M. R. Mousavi, T. Neele and M. Sammartino,
\emph{Compositional Active Learning of Synchronizing Systems Through
Automated Alphabet Refinement}, CONCUR 2025.
\url{https://drops.dagstuhl.de/storage/00lipics/lipics-vol348-concur2025/LIPIcs.CONCUR.2025.20/LIPIcs.CONCUR.2025.20.pdf}.
\bibitem{MinimalRealization} R. Rosebrugh, N. Sabadini and R. F. C. Walters,
\emph{Minimal Realization in Bicategories of Automata},
Mathematical Structures in Computer Science (1998).
\url{https://doi.org/10.1017/S0960129597002454}.
\bibitem{Common} Previous project manuscript,
\emph{One operational construction for residual automata, formal concepts,
and topology}, 4 November 2026.
Recovered source \texttt{common\_generation\_note.tex}, included in the archive.
\bibitem{Wiring} Previous project manuscript,
\emph{Contextual Reduction and Wiring for Open Learners}, 5 November 2026.
Recovered source \texttt{contextual\_reduction\_wiring.tex}, included in the archive.
\bibitem{Adaptive} Previous project manuscript,
\emph{Evidence Guided Native Experiment Formation}, 5 November 2026.
Recovered source \texttt{adaptive\_experiment\_formation.tex}, included in the archive.
\end{thebibliography}
\end{document}
